% ============================================================================
%  EHD推進 計算シミュレーション実験計画書
% ============================================================================
\documentclass[12pt,a4paper]{ltjsarticle}

\usepackage{luatexja-fontspec}
\setmainjfont{HaranoAjiMincho-Regular}[BoldFont=HaranoAjiMincho-Bold]
\setsansjfont{HaranoAjiGothic-Medium}[BoldFont=HaranoAjiGothic-Bold]

\usepackage{amsmath,amssymb}
\usepackage{graphicx}
\usepackage{hyperref}
\usepackage[margin=2cm]{geometry}
\usepackage{booktabs}
\usepackage{enumitem}
\usepackage{xcolor}
\usepackage{tcolorbox}
\usepackage{fancyhdr}
\usepackage{tikz}
\usetikzlibrary{arrows.meta,shapes.geometric,positioning,calc}

\tcbuselibrary{skins,breakable}

\definecolor{tier1}{RGB}{52,152,219}
\definecolor{tier2}{RGB}{230,126,34}
\definecolor{tier3}{RGB}{142,68,173}
\definecolor{info}{RGB}{0,123,255}
\definecolor{success}{RGB}{40,167,69}
\definecolor{warning}{RGB}{255,193,7}
\definecolor{cloud}{RGB}{0,150,136}
\definecolor{lightgray}{RGB}{245,245,245}

\newtcolorbox{infobox}[1][]{
  colback=info!8, colframe=info!70,
  fonttitle=\bfseries\sffamily, title={#1}, breakable
}
\newtcolorbox{successbox}[1][]{
  colback=success!8, colframe=success!70,
  fonttitle=\bfseries\sffamily, title={#1}, breakable
}
\newtcolorbox{warningbox}[1][]{
  colback=warning!15, colframe=warning!80!black,
  fonttitle=\bfseries\sffamily, title={#1}, breakable
}
\newtcolorbox{tierbox}[2][]{
  colback=#2!8, colframe=#2!80,
  fonttitle=\bfseries\sffamily\Large, title={#1}, breakable,
  before upper={\parindent=0pt}
}

\pagestyle{fancy}
\fancyhf{}
\fancyhead[L]{\sffamily\small EHD推進 計算シミュレーション実験計画書}
\fancyhead[R]{\sffamily\small v1.0}
\fancyfoot[C]{\thepage}

\begin{document}

% ============================================================================
%  TITLE
% ============================================================================
\begin{titlepage}
\centering
\vspace*{2cm}

{\Huge\sffamily\bfseries
EHD推進の\\[8pt]
計算シミュレーション\\[8pt]
実験計画書
}

\vspace{1cm}

{\Large\sffamily
--- プラズマを仮想空間で再現し、制御する ---
}

\vspace{2cm}

\begin{tikzpicture}[scale=0.9]
  % Grid/mesh representation
  \foreach \x in {0,0.5,...,6} {
    \draw[gray!30, thin] (\x,0) -- (\x,4);
  }
  \foreach \y in {0,0.5,...,4} {
    \draw[gray!30, thin] (0,\y) -- (6,\y);
  }

  % Density contour (fake)
  \fill[blue!30, opacity=0.5, rounded corners=3pt] (1,1) -- (2,0.5) -- (4,0.8) -- (5,1.5) -- (4.5,3) -- (2.5,3.5) -- (1,2.5) -- cycle;
  \fill[red!40, opacity=0.5, rounded corners=3pt] (2,1.2) -- (3,1) -- (4,1.5) -- (3.5,2.5) -- (2.5,2.8) -- cycle;
  \fill[yellow!50, opacity=0.6, rounded corners=2pt] (2.5,1.5) -- (3.2,1.3) -- (3.5,2) -- (3,2.3) -- cycle;

  % Labels
  \node[font=\sffamily\small, white] at (3,1.8) {\textbf{プラズマ}};
  \node[font=\sffamily\tiny, blue!70!black] at (1.5,3) {$n_i$};
  \node[font=\sffamily\tiny, red!70!black] at (4.2,2.8) {$\mathbf{E}$};

  % Electrodes
  \draw[very thick, blue] (0.2,4) -- (5.8,4);
  \draw[very thick, orange!80!black] (0.2,0) -- (5.8,0);
  \node[above, font=\sffamily\small, blue] at (3,4) {エミッター};
  \node[below, font=\sffamily\small, orange!80!black] at (3,0) {コレクター};

  % GPU icon
  \node[draw, thick, fill=success!20, rounded corners=4pt,
    font=\sffamily\small, minimum width=1.5cm] at (8,2) {GPU};
  \draw[-{Stealth}, thick, gray] (6.3,2) -- (7.2,2);

  % Monitor
  \draw[thick, rounded corners=2pt] (9.5,0.5) rectangle (12.5,3.5);
  \draw[thick] (10.5,0.5) -- (10.5,0) -- (11.5,0) -- (11.5,0.5);
  \node[font=\sffamily\tiny] at (11,2) {リアルタイム};
  \node[font=\sffamily\tiny] at (11,1.5) {可視化};
\end{tikzpicture}

\vspace{2cm}

{\large\sffamily
対象：プログラミング経験のある方\\
予算：無料（クラウド）〜 約60万円（自作PC）\\
物理実験の危険なし --- 全てPC上で完結
}

\vspace{1cm}
{\sffamily\small 2026年3月}
\end{titlepage}

% ============================================================================
%  OVERVIEW
% ============================================================================
\section{この計画書でやること}

\subsection{目的}

論文で提案したEHD推進のプラズマ制御理論を、
\textbf{物理実験なし}で計算機上に再現し検証します。

\begin{center}
\begin{tikzpicture}[
  stepbox/.style={draw, rounded corners=6pt, minimum width=3.8cm,
    minimum height=1.6cm, font=\sffamily\small, align=center, thick},
  arr/.style={-{Stealth[length=5pt]}, thick, gray}
]
  \node[stepbox, fill=tier1!15, draw=tier1] (s1) at (0,0)
    {\textbf{Step 1}\\2Dコロナ放電\\シミュレーション};
  \node[stepbox, fill=tier2!15, draw=tier2] (s2) at (5,0)
    {\textbf{Step 2}\\不安定性解析\\成長率の計算};
  \node[stepbox, fill=tier3!15, draw=tier3] (s3) at (10,0)
    {\textbf{Step 3}\\フィードバック\\制御の実装};
  \draw[arr] (s1) -- (s2);
  \draw[arr] (s2) -- (s3);
\end{tikzpicture}
\end{center}

\subsection{Isaac Simは使えるか？}

\begin{warningbox}[結論：プラズマシミュレーションにはIsaac Simは不向き]
NVIDIA Isaac Simはロボティクス向けであり、プラズマ物理（電離、
イオンドリフト、コロナ放電）の第一原理ソルバーを持ちません。
PhysXは剛体/関節力学、Flowは煙・炎のビジュアル用簡易流体であり、
ナビエ-ストークス--マクスウェル連成方程式は解けません。

\textbf{代わりに使うべきツール：}
\begin{itemize}[leftmargin=2em]
  \item \textbf{ZAPDOS}（無料）--- 低温プラズマ流体の本命
  \item \textbf{WarpX}（無料）--- GPU加速PICコード
  \item \textbf{OpenFOAM + interEHDFoam}（無料）--- EHD流体
  \item \textbf{COMSOL Plasma Module}（有料）--- 最も完成度の高い商用ツール
\end{itemize}
\end{warningbox}

% ============================================================================
%  SOFTWARE
% ============================================================================
\newpage
\section{シミュレーションソフトウェア選択ガイド}

\subsection{用途別フローチャート}

\begin{center}
\begin{tikzpicture}[
  node distance=1.2cm,
  q/.style={draw, diamond, aspect=2.5, inner sep=2pt,
    font=\sffamily\small, align=center, thick, fill=lightgray},
  a/.style={draw, rounded corners=4pt, font=\sffamily\small,
    align=center, thick, minimum width=3cm},
  arr/.style={-{Stealth[length=4pt]}, thick}
]
  \node[q] (q1) {やりたいこと};

  \node[a, fill=tier1!20, below left=1.5cm and 2cm of q1] (a1)
    {\textbf{ZAPDOS}\\(無料・FEM)};
  \node[a, fill=tier2!20, below=1.5cm of q1] (a2)
    {\textbf{WarpX}\\(無料・GPU PIC)};
  \node[a, fill=tier3!20, below right=1.5cm and 2cm of q1] (a3)
    {\textbf{OpenFOAM}\\(無料・CFD)};

  \draw[arr] (q1) -- node[left, font=\sffamily\tiny, align=right]
    {コロナ放電\\プラズマ流体} (a1);
  \draw[arr] (q1) -- node[right, font=\sffamily\tiny, align=center]
    {粒子運動論\\PIC計算} (a2);
  \draw[arr] (q1) -- node[right, font=\sffamily\tiny, align=left]
    {EHD流体\\CFD連成} (a3);
\end{tikzpicture}
\end{center}

\subsection{ソフトウェア比較表}

\begin{table}[h]
\centering\sffamily\small
\caption{プラズマ/EHDシミュレーションソフトウェア一覧}
\begin{tabular}{@{}p{2cm}p{1.5cm}p{2.5cm}p{2cm}p{3cm}@{}}
\toprule
\textbf{ソフト} & \textbf{費用} & \textbf{手法} & \textbf{GPU対応} & \textbf{得意分野} \\
\midrule
\textbf{ZAPDOS} & 無料 & FEM（MOOSE） & CPU中心 & コロナ放電、DBD、低温プラズマ流体 \\
\textbf{WarpX} & 無料 & PIC & CUDA/HIP & 粒子運動論、大規模プラズマ \\
\textbf{OpenFOAM} & 無料 & 有限体積法 & CPU中心 & EHD二相流、CFD \\
\textbf{FEniCS} & 無料 & FEM & CPU & カスタムPDE、ドリフト拡散 \\
\textbf{FiPy} & 無料 & 有限体積(Python) & CPU & プロトタイプ実装 \\
\textbf{COMSOL} & 有料 & FEM & CPU & 最も完成度の高い商用ツール \\
\textbf{ANSYS} & 有料 & FEM/FVM & CPU+GPU & DBDアクチュエータ \\
\bottomrule
\end{tabular}
\end{table}

\subsection{各ソフトの詳細}

\subsubsection{ZAPDOS --- 本計画の中核ツール（推奨）}

\begin{successbox}[最も推奨：EHDプラズマシミュレーションの最適解]
\begin{itemize}[leftmargin=2em]
  \item 米国Idaho国立研究所（INL）開発のオープンソース低温プラズマソルバー
  \item MOOSEフレームワーク上で動作（自動微分、並列FEM、メッシュ適応）
  \item \textbf{CRANE}（プラズマ化学）と\textbf{ELK}（電磁場）と連成可能
  \item 電子・イオンのドリフト-拡散-反応方程式を直接解ける
  \item 論文の式(5)--(6)をそのまま実装可能
  \item GitHub: \texttt{github.com/shannon-lab/zapdos}
  \item OS: Linux推奨、macOS/Windows(WSL)も可
\end{itemize}
\end{successbox}

\subsubsection{WarpX --- GPU加速PICコード}

\begin{infobox}[粒子レベルのプラズマ物理を見たい場合]
\begin{itemize}[leftmargin=2em]
  \item 2022年ACM Gordon Bell賞受賞の最先端PICコード
  \item CUDA（NVIDIA）、HIP（AMD）、DPC++（Intel）全対応
  \item 単一GPU〜Exascaleスパコンまでスケール
  \item 最低VRAM: 8GB（小規模2D）、推奨: 24GB以上
  \item GitHub: \texttt{github.com/BLAST-WarpX/warpx}
\end{itemize}
\end{infobox}

\subsubsection{OpenFOAM + interEHDFoam --- EHD流体}

EHDの体積力項（$\mathbf{f}_{\mathrm{EHD}} = \rho_q \mathbf{E}$）を
含むCFDシミュレーション。OpenFOAM自体は広く使われており
日本語の情報も豊富。

\begin{itemize}[leftmargin=2em]
  \item \texttt{github.com/ntofighi/interEHDFoam}（二相EHD流れ）
  \item \texttt{github.com/silviomcandido/interIsoFoamEHD}（高解像度版）
  \item 日本語資料：OpenFOAM JP Wiki、CFD-Onlineフォーラム
\end{itemize}

\subsubsection{COMSOL Plasma Module --- 商用最高峰}

\begin{itemize}[leftmargin=2em]
  \item コロナ放電、DBD、グロー放電をGUIで構築可能
  \item アカデミック年間ライセンス: 約\$1,495（約22万円）
  \item 個人には高価だが、TSUBAME経由なら月4ポイント（約250円）で利用可
  \item 日本法人あり：comsol.com/japan
\end{itemize}

% ============================================================================
%  HARDWARE
% ============================================================================
\newpage
\section{ハードウェア：PC構成ガイド}

\subsection{3段階のPC構成}

\begin{tierbox}[Tier 1：エントリー --- 2Dシミュレーション入門]{tier1}
\textbf{予算：17〜25万円} \quad
\textbf{対象：}ZAPDOSで2Dコロナ放電を試す

\medskip\centering
\centering\sffamily\small
\begin{tabular}{@{}llr@{}}
\toprule
\textbf{パーツ} & \textbf{推奨品} & \textbf{価格目安} \\
\midrule
CPU & AMD Ryzen 9 7900X (12C/24T) & 40,000〜55,000円 \\
GPU & NVIDIA RTX 4070 Ti Super (16GB) & 70,000〜100,000円 \\
RAM & DDR5-5600 32GB (16$\times$2) & 20,000〜35,000円 \\
SSD & NVMe 1TB PCIe 4.0 & 10,000〜15,000円 \\
マザーボード & AMD AM5 ATX & 20,000〜30,000円 \\
電源 & 850W 80+ Gold & 12,000〜18,000円 \\
\midrule
\textbf{合計} & & \textbf{17〜25万円} \\
\bottomrule
\end{tabular}
\medskip

\textbf{できること：}
ZAPDOS 1D/2D、FEniCS/FiPyカスタム実装、WarpX小規模2D PIC（\textasciitilde 10M粒子）
\end{tierbox}

\vspace{0.5cm}

\begin{tierbox}[Tier 2：推奨 --- 2D軸対称＋フィードバック制御]{tier2}
\textbf{予算：48〜63万円} \quad
\textbf{対象：}論文の理論を本格的に検証する

\medskip\centering
\centering\sffamily\small
\begin{tabular}{@{}llr@{}}
\toprule
\textbf{パーツ} & \textbf{推奨品} & \textbf{価格目安} \\
\midrule
CPU & AMD Ryzen 9 9950X (16C/32T) & 75,000〜90,000円 \\
GPU & \textbf{NVIDIA RTX 4090 (24GB)} & 290,000〜350,000円 \\
RAM & DDR5-6000 64GB (32$\times$2) & 40,000〜70,000円 \\
SSD & NVMe 2TB + HDD 4TB & 20,000〜30,000円 \\
マザーボード & AMD X670E ATX & 30,000〜45,000円 \\
電源 & 1000W 80+ Platinum & 18,000〜25,000円 \\
CPUクーラー & 360mm AIO水冷 & 12,000〜20,000円 \\
\midrule
\textbf{合計} & & \textbf{48〜63万円} \\
\bottomrule
\end{tabular}
\medskip

\textbf{できること：}
WarpX 2D〜小規模3D PIC（\textasciitilde 100M粒子）、ZAPDOS+ELK連成、
フィードバック制御の数値実験、COMSOL（TSUBAME経由）
\end{tierbox}

\vspace{0.5cm}

\begin{tierbox}[Tier 3：ハイエンド --- フル3Dシミュレーション]{tier3}
\textbf{予算：90〜200万円} \quad
\textbf{対象：}3D完全シミュレーション

\medskip\centering
\centering\sffamily\small
\begin{tabular}{@{}llr@{}}
\toprule
\textbf{パーツ} & \textbf{推奨品} & \textbf{価格目安} \\
\midrule
CPU & AMD Threadripper PRO 7965WX (24C) & 250,000〜400,000円 \\
GPU & \textbf{RTX 5090 (32GB)} or RTX 6000 Ada (48GB) & 400,000〜800,000円 \\
RAM & DDR5-6400 ECC 128GB & 63,000〜100,000円 \\
SSD & NVMe 4TB PCIe 5.0 + 2TB secondary & 60,000〜100,000円 \\
マザーボード & TRX50 (Threadripper対応) & 80,000〜150,000円 \\
電源 & 1600W 80+ Titanium & 35,000〜55,000円 \\
CPUクーラー & 420mm AIO or カスタム水冷 & 20,000〜50,000円 \\
\midrule
\textbf{合計} & & \textbf{90〜200万円} \\
\bottomrule
\end{tabular}
\medskip

\textbf{できること：}
WarpX/PIConGPU完全3D PIC（\textasciitilde 1B粒子）、ZAPDOS+CRANE+ELK 3D連成
\end{tierbox}

\subsection{パーツ購入先（日本国内）}

\begin{table}[h]
\centering\sffamily\small
\caption{PCパーツ購入先一覧}
\begin{tabular}{@{}lll@{}}
\toprule
\textbf{店舗} & \textbf{特徴} & \textbf{URL} \\
\midrule
ツクモ (Tsukumo) & 自作PC専門、秋葉原に実店舗 & shop.tsukumo.co.jp \\
ドスパラ (Dospara) & BTO対応、全国に店舗 & dospara.co.jp \\
パソコン工房 & BTO対応、サポート充実 & pc-koubou.jp \\
Amazon.co.jp & 品揃え豊富、即日配送 & amazon.co.jp \\
ソフマップ & 中古パーツも扱う & sofmap.com \\
\bottomrule
\end{tabular}
\end{table}

\begin{warningbox}[2026年3月現在の注意]
\begin{itemize}[leftmargin=2em]
  \item \textbf{DDR5 RAM}：AI需要の急増により世界的に不足中。
        ツクモ等で購入制限・価格高騰が発生。早めの確保を推奨
  \item \textbf{RTX 5090}：2025年1月発売後、抽選販売が続いた。
        2025年6月以降398,000円以下の製品も出始めたが依然入手困難。
        \textbf{RTX 4090の中古（20万円〜）}も有力な選択肢
\end{itemize}
\end{warningbox}

% ============================================================================
%  CLOUD
% ============================================================================
\newpage
\section{クラウド計算：PCを買わない選択肢}

\begin{infobox}[PCを自作せずにプラズマシミュレーションを始める方法]
高性能GPUを持っていなくても、クラウドサービスを使えば
\textbf{今日から}シミュレーションを始められます。
\end{infobox}

\subsection{サービス比較}

\begin{table}[h]
\centering\sffamily\small
\caption{クラウドGPU計算サービス比較}
\begin{tabular}{@{}p{2.5cm}p{2cm}p{2cm}p{2cm}p{3cm}@{}}
\toprule
\textbf{サービス} & \textbf{GPU} & \textbf{月額/時間} & \textbf{最低費用} & \textbf{備考} \\
\midrule
\textbf{Google Colab Pro} & T4/V100/A100 & 約1,500円/月 & 1,500円 & 最も手軽。ZAPDOS/FiPy向き \\
\textbf{Colab Pro+} & A100優先 & 約7,500円/月 & 7,500円 & 長時間連続実行可能 \\
\textbf{AWS EC2 g5} & A10G (24GB) & 約150円/時 & --- & WarpX小〜中規模 \\
\textbf{AWS EC2 p4d} & A100$\times$8 & 約5,000円/時 & --- & 大規模3D PIC \\
\textbf{ABCI 3.0} & H200$\times$8 & 3,520円/時 & 22万円 & 国内最強。研究機関限定 \\
\textbf{TSUBAME 4.0} & H100相当 & 極めて安価 & 学内無料 & 東京科学大学所属者 \\
\bottomrule
\end{tabular}
\end{table}

\subsection{おすすめルート}

\begin{center}
\begin{tikzpicture}[
  box/.style={draw, rounded corners=6pt, minimum width=4.5cm,
    minimum height=1.8cm, font=\sffamily\small, align=center, thick},
  arr/.style={-{Stealth[length=5pt]}, thick, gray}
]
  \node[box, fill=success!15, draw=success] (free) at (0,0)
    {\textbf{無料〜月1,500円}\\Google Colab Pro\\+ ZAPDOS/FiPy};
  \node[box, fill=tier2!15, draw=tier2] (mid) at (6,0)
    {\textbf{月5千〜5万円}\\AWS スポットGPU\\+ WarpX};
  \node[box, fill=tier3!15, draw=tier3] (high) at (12,0)
    {\textbf{大学所属者}\\ABCI / TSUBAME\\+ COMSOL};

  \node[below, font=\sffamily\tiny] at (free.south) {誰でも今すぐ};
  \node[below, font=\sffamily\tiny] at (mid.south) {本格計算};
  \node[below, font=\sffamily\tiny] at (high.south) {最高性能};

  \draw[arr] (free) -- (mid);
  \draw[arr] (mid) -- (high);
\end{tikzpicture}
\end{center}

\begin{successbox}[大学・研究機関所属者は圧倒的に有利]
\begin{itemize}[leftmargin=2em]
  \item \textbf{TSUBAME 4.0}（東京科学大学）：学生は実質無料でH100クラスGPU利用可。
        COMSOLも月250円相当で利用可能
  \item \textbf{ABCI 3.0}（産総研）：H200$\times$6,128基。
        産官学の研究機関が対象。最低22万円から
  \item \textbf{mdx}（9国立大連携）：IaaSクラウド。所属大学経由でアクセス
  \item いずれも商用クラウドより\textbf{桁違いに安い}
\end{itemize}
\end{successbox}

% ============================================================================
%  SIMULATION PLAN
% ============================================================================
\newpage
\section{シミュレーション実験の手順}

\subsection{Step 1：2Dコロナ放電シミュレーション}

\begin{tierbox}[Step 1 --- 基本的なEHDプラズマの再現]{tier1}
\textbf{ツール：}ZAPDOS（推奨）or FiPy \quad
\textbf{所要時間：}1〜2週間 \quad
\textbf{計算資源：}Colab Pro or Tier 1 PC
\end{tierbox}

\subsubsection{セットアップ}

\begin{enumerate}[leftmargin=2em]
  \item MOOSEフレームワークのインストール（mooseframework.inl.gov）
  \item ZAPDOSのビルド（GitHub: shannon-lab/zapdos）
  \item チュートリアル問題の実行（1Dグロー放電）
\end{enumerate}

\subsubsection{EHDコロナ放電モデルの構築}

\begin{itemize}[leftmargin=2em]
  \item \textbf{計算領域：}2D軸対称、電極間ギャップ$d = 1$\,cm
  \item \textbf{エミッター：}曲率半径$r = 0.1$\,mm（高電場を模擬）
  \item \textbf{コレクター：}平板電極
  \item \textbf{ガス：}大気圧空気（$N_2$/O$_2$ = 79/21）
  \item \textbf{印加電圧：}$V = 10$--$30$\,kV（パラメトリックスイープ）
  \item \textbf{種：}$e^-$, $\mathrm{N}_2^+$, $\mathrm{O}_2^-$（3種モデル）
\end{itemize}

\subsubsection{解くべき方程式（論文との対応）}

論文の式(5)--(6)に直接対応する種輸送方程式：
\begin{align}
  \frac{\partial n_i}{\partial t}
  + \nabla \cdot (n_i \mathbf{v} + n_i \mu_i \mathbf{E} - D_i \nabla n_i)
  &= \alpha_T |\Gamma_e| - \beta_{\mathrm{rec}} n_i n_e \\
  \frac{\partial n_e}{\partial t}
  + \nabla \cdot (n_e \mathbf{v} - n_e \mu_e \mathbf{E} - D_e \nabla n_e)
  &= \alpha_T |\Gamma_e| - \beta_{\mathrm{rec}} n_i n_e
\end{align}
ポアソン方程式（式4の静電近似）：
\begin{equation}
  \nabla^2 \phi = -\frac{e(n_i - n_e)}{\varepsilon_0}
\end{equation}

\subsubsection{検証項目}

\begin{table}[h]
\centering\sffamily\small
\begin{tabular}{@{}p{4cm}p{4cm}p{3.5cm}@{}}
\toprule
\textbf{測定量} & \textbf{期待される結果} & \textbf{論文との対応} \\
\midrule
コロナ開始電圧 & 空気中10〜15\,kV & 既知の実験値と比較 \\
電流-電圧特性 & $I \propto V(V - V_{\mathrm{onset}})$ & Townsendの関係 \\
イオン密度分布 & ドリフト領域で$n_i \sim 10^{18}\,\mathrm{m}^{-3}$ & 論文の仮定値 \\
EHD体積力分布 & $f = \rho_q E$ & 式(1) \\
\bottomrule
\end{tabular}
\end{table}

\subsection{Step 2：不安定性解析}

\begin{tierbox}[Step 2 --- 成長率の計算と論文予測の検証]{tier2}
\textbf{ツール：}ZAPDOS + Python（後処理） \quad
\textbf{所要時間：}2〜4週間 \quad
\textbf{計算資源：}Tier 1〜2 PC
\end{tierbox}

\subsubsection{線形安定性解析}

\begin{enumerate}[leftmargin=2em]
  \item Step 1で得た定常解を基本状態$\mathbf{U}_0$とする
  \item 小さな摂動を加えて時間発展を追跡
  \item 摂動の指数関数的成長率$\gamma$を抽出
  \item 論文の予測値と比較：
    \begin{itemize}
      \item $\gamma_{\mathrm{ion}} \lesssim 10^9\,\mathrm{s}^{-1}$（式15）
      \item $\gamma_{\mathrm{EHD}} \lesssim 5 \times 10^7\,\mathrm{s}^{-1}$（式14）
    \end{itemize}
\end{enumerate}

\begin{infobox}[論文の予測が上界推定であることの検証]
論文では$\gamma_{\mathrm{ion}} \sim 10^9\,\mathrm{s}^{-1}$と推定していますが、
これは電子拡散・再結合・衝突減衰を無視した上界です。
シミュレーションではこれらの効果が自動的に含まれるため、
\textbf{実効的な成長率は論文の予測より小さくなるはず}です。
この差の定量化が重要な検証ポイントです。
\end{infobox}

\subsection{Step 3：フィードバック制御の実装}

\begin{tierbox}[Step 3 --- 制御による安定範囲拡大の検証]{tier3}
\textbf{ツール：}ZAPDOS + Python制御ループ \quad
\textbf{所要時間：}1〜3ヶ月 \quad
\textbf{計算資源：}Tier 2 PC以上
\end{tierbox}

\subsubsection{制御ループの実装}

\begin{enumerate}[leftmargin=2em]
  \item コレクター電極を4セグメントに分割（計算メッシュ上で）
  \item 各セグメントの電流密度を「センサー」として読み取り
  \item PythonスクリプトからZAPDOSを呼び出す外側ループ：
    \begin{itemize}
      \item 現在の電流分布を取得
      \item 摂動を検出（基準値からの偏差）
      \item 電極電圧を微調整（PID制御 or LQR制御）
      \item 次のタイムステップを実行
    \end{itemize}
  \item \textbf{制御なし}と\textbf{制御あり}でアーク遷移電圧を比較
\end{enumerate}

\subsubsection{検証すべき論文の予測}

\begin{table}[h]
\centering\sffamily\small
\begin{tabular}{@{}p{4.5cm}p{4cm}p{3cm}@{}}
\toprule
\textbf{論文の予測} & \textbf{シミュレーションで確認} & \textbf{対応箇所} \\
\midrule
制御で安定範囲が拡大 & $V_{\mathrm{arc}}^{\mathrm{ctrl}} > V_{\mathrm{arc}}^{\mathrm{open}}$ & 命題1 \\
$\mathrm{Re}_{\max} \propto f_s^{2/3}$ & 制御帯域幅を変えてスケーリング確認 & 命題5 \\
エントロピーがドリフト領域から電極へ移動 & $\Omega_d$と$\Omega_c$のエントロピー生成比較 & 命題1(ii) \\
4セグメントで十分 & セグメント数 vs 安定化効果 & 系4 \\
\bottomrule
\end{tabular}
\end{table}

% ============================================================================
%  PYTHON STARTER
% ============================================================================
\newpage
\section{すぐ始められるPythonコード}

ZAPDOSのセットアップに時間がかかる場合、まず\textbf{FiPy}（Python有限体積法ライブラリ）で
簡易版を動かすことをお勧めします。

\begin{successbox}[5分で始める：pip install fipy]
\texttt{pip install fipy numpy matplotlib}

FiPyで1Dドリフト-拡散方程式を解き、イオン密度の時間発展を可視化できます。
コロナ放電の基本物理を理解するための第一歩として最適です。

\vspace{0.5cm}
学習リソース：
\begin{itemize}[leftmargin=2em]
  \item FiPy公式ドキュメント: \texttt{www.ctcms.nist.gov/fipy/}
  \item PlasmaPy（プラズマ物理ユーティリティ）: \texttt{www.plasmapy.org}
  \item Plasma Simulations by Example: \texttt{www.plasmasimulationsbyexample.com}
  \item PICチュートリアル（ゼロから実装）: \texttt{www.particleincell.com}
\end{itemize}
\end{successbox}

% ============================================================================
%  SUMMARY
% ============================================================================
\section{全体まとめ}

\begin{table}[h]
\centering\sffamily\small
\caption{予算・目的別クイックリファレンス}
\begin{tabular}{@{}p{3cm}p{3.5cm}p{5cm}@{}}
\toprule
\textbf{予算} & \textbf{構成} & \textbf{できること} \\
\midrule
無料〜月1,500円 & Colab Pro + FiPy & 1D/2Dプロトタイプ、基本物理の理解 \\
17〜25万円 & Tier 1 自作PC & ZAPDOS 2D、WarpX小規模2D \\
48〜63万円 & Tier 2 自作PC & 本格2D軸対称、フィードバック制御実験 \\
90〜200万円 & Tier 3 自作PC & フル3Dシミュレーション \\
大学所属 & TSUBAME/ABCI & 最高性能 + COMSOL（実質無料〜） \\
\bottomrule
\end{tabular}
\end{table}

\begin{infobox}[最初の一歩]
\begin{enumerate}[leftmargin=2em]
  \item \textbf{今日}：Google Colabで \texttt{pip install fipy} を実行し、1Dドリフト拡散を解く
  \item \textbf{今週}：ZAPDOSをインストールし、チュートリアル（1Dグロー放電）を動かす
  \item \textbf{今月}：2Dコロナ放電モデルを構築し、電流-電圧特性を計算する
  \item \textbf{3ヶ月後}：フィードバック制御を実装し、安定範囲の拡大を検証する
\end{enumerate}
物理実験の危険なし。必要なのはPC（またはクラウドアカウント）と好奇心だけです。
\end{infobox}

\end{document}
